% Fragmento autónomo para el borrador del artículo II, 8 de septiembre de 2026.
% Procedencia: RESULTADO_RECUPERADO y FORMALIZACION_REUNIDA de los propietarios
% B07/c41_ckm.tex y C02/74_mezcla_contenido_exclusivo_post_rev2.tex del
% BORRADOR_02; integral 2249, caps. 49, 92.6, 92.24 y 93.10--93.11.
% Los originales y su soporte demostrativo íntegro no se sustituyen ni editan.
% Claves bibliográficas utilizadas: ckm, mezcla, pmnsespectral, cpt.
% Necesita: amsmath, amssymb, amsthm; theorem/proposition/definition/remark.

\section{Ángulos de mezcla CKM, fase CP y transporte leptónico}
\label{sec:ii-propagacion-angular}

\subsection{Antecedente común y separación de los dos transportes}

El fundamento común del artículo produce el estado enriquecido mediante
las dos evaluaciones de APP, la selección orientada del TRIT y la
composición efectiva de selección, transporte, acarreo, memoria y retorno
del TPK. Las representaciones angulares no reemplazan ese estado: son lectores
de su registro. En particular, la prolongación
\(w_6\to w_{12}\to w_{18}\to w_{24}\to w_{30}\to R_{36}\to G_9\)
conserva la procedencia de las lecturas numéricas, y la holonomía
\(\Gamma_9\) actúa antes de la representación reducida \(K_{\rm ph}\).
El retorno de fase no reinicia el registro.

De ese mismo antecedente se reciben aquí dos clases de objetos. La primera
contiene el ángulo directo, el contraángulo y el defecto torsional
\(\Delta_4\), ya construidos en sus lectores propios. La segunda contiene
las extensiones supervivientes, sus rutas observables, el registro
dodecafásico y los operadores de las fibras materiales. La cadena
extensión--ruta--registro--firma--carácter--torres conserva el atlas
\(81/56/13/324\) y distingue la coordenatización nula de la rama positiva.
Los ángulos de mezcla no eligen retrospectivamente esas rutas ni generan
los autovalores de masa.

El transporte quark utiliza una coordenatización angular racional y su realización
unitaria compleja. El transporte leptónico utiliza los marcos espectrales
de las fibras cargada y neutra. Ambos proceden de la estructura discreta
común, pero sus dominios, sus lectores y sus cambios de base son distintos.
El funcional de Barbero--Immirzi y la mezcla tampoco se definen uno por el
otro: reciben representaciones tipadas del antecedente angular y de incidencia.
La relación entre ambos se expresa mediante esos mapas, no mediante una
fórmula retrospectiva que use un valor de Barbero como selector de mezcla.

\begin{remark}[Reglas sectoriales y evaluación del lector]
Las tres reglas sectoriales y su evaluación se desarrollan a continuación
en la sección~\ref{sec:ii09-construccion-sectorial}. El teorema
\ref{thm:ii09-matriz-sectorial} obtiene las columnas como imágenes de la
base angular. Las pruebas posteriores establecen la inversión de la coordenatización,
el transporte de hoja, la métrica y la realización compleja.
La elección de las reglas y la evaluación de sus coeficientes son dos
cuestiones matemáticas distintas: el sistema de partida se declara en la
Definición~\ref{def:ii09-sistema-sectorial}, y la prueba del teorema
utiliza ese sistema completo. La composición previa que selecciona esas
reglas a partir de la incidencia requiere su demostración propia.
\end{remark}

% Aportación separada, 9 de septiembre de 2026.
% RESULTADO_RECUPERADO: construcción sectorial del propietario 08_ckm_cp.tex.
% FORMALIZACION_REUNIDA: descomposición de paridad y compatibilidad representacional.
% Insertar antes de la carta racional CKM; conservar el desarrollo posterior.
% La procedencia y el alcance exacto están en 08b_procedencia_angular.md.

\subsection{Incidencia sectorial y construcción del lector angular}
\label{sec:ii09-construccion-sectorial}

La construcción angular utiliza las representaciones de acción y cierre
obtenidas previamente desde APP--TRIT--TPK. El ángulo directo, el
contraángulo y el defecto torsional conservan sus dominios propios.
El lector quark combina estas tres coordenadas mediante la incidencia
sectorial; su evaluación precede a las amplitudes complejas y al
contraste metrológico. La inclusión hexádico--octádica interviene en
el tercer sector a través de la diferencia de sus cardinales.

% ARQUITECTURA_AUTORAL_PREEXISTENTE: semicuadrado bidireccional de N42.
% FORMALIZACION_NUEVA: grupoide de pares y ponderacion por estabilizadores.
\subsubsection{Incidencia marcada y normalización bidireccional}
\label{sec:ii09-normalizacion-bidireccional}

Los invariantes del lector conservan la naturaleza del objeto del que
proceden. La intersección marcada $B_K=H_e\cap P_\pi$ tiene cardinal
$b=4$; el entero $q=7$ es la coordenada del pivote $p_\varphi$ en la
numeración orientada. El orden de determinación del diseño de Steiner
es $p=5$. Los soportes hexádico y octádico tienen cardinales $h_6=6$
y $o_8=8$, y el alfabeto de orientación trítica tiene cardinal $t=3$.
Un cardinal de soporte y una coordenada de pivote se utilizan mediante
sus aplicaciones respectivas, aun cuando ambos se publiquen como enteros.

La normalización del autocruce pentagonal admite la siguiente
realización explícita. Sea $P$ el soporte de cinco posiciones del cierre
y sea $X=P\times P$ su espacio de pares ordenados. La inversión
bidireccional sobre este espacio es
\[
 \jmath:X\longrightarrow X,\qquad \jmath(x,y)=(y,x),
 \qquad\jmath^2=I.
\]
Se conserva la información de estabilizador de cada órbita. Para una
órbita $[z]$ se define su peso como
$1/|\operatorname{Stab}_{C_2}(z)|$. La suma de esos pesos es la
cardinalidad del grupoide de acción $X/\!/C_2$.

\begin{proposition}[Medida orbital del autocruce bidireccional]
\label{prop:ii09-semicuadrado-orbital}
Para un soporte finito de cardinal $p$, la inversión de pares ordenados
determina
\[
 |(P\times P)/\!/C_2|
 =\sum_{[z]\in(P\times P)/C_2}
      \frac1{|\operatorname{Stab}_{C_2}(z)|}
 =\frac{p^2}{2}.
\]
Para $p=5$, el autocruce tiene diez órbitas libres y cinco órbitas
fijas, y su medida orbital es $10+5/2=25/2$.
\end{proposition}
\begin{proof}
Cada par diagonal $(x,x)$ forma una órbita fija, cuyo estabilizador
tiene orden dos; hay $p$ de ellas. Los $p(p-1)$ pares restantes se
agrupan de dos en dos, con estabilizador trivial. Su contribución es
por tanto $p(p-1)/2$, y la de la diagonal es $p/2$. La suma es
$p^2/2$. Equivalentemente, la identidad órbita--estabilizador da
$1/|\operatorname{Stab}(z)|=|[z]|/2$; sumar sobre todas las órbitas
produce $|P\times P|/2$.
\end{proof}

La medida conserva explícitamente la contribución de los puntos fijos.
El conjunto ordinario de órbitas tiene cardinal $p(p+1)/2$, que vale
quince en el cierre pentagonal. La medida orbital y ese cardinal son
dos lecturas determinadas de la misma inversión. La normalización
bidireccional del lector utiliza la primera: las diez órbitas libres
pesan una unidad y las cinco diagonales pesan media unidad.

Esta realización precisa el factor $p^2/2$ empleado en el sistema
sectorial siguiente. La asignación de esa contribución al sector $23$,
su orientación negativa y las otras dos reglas de incidencia permanecen
especificadas por el sistema sectorial completo. La prueba establece la
normalización del autocruce y preserva diferenciada la selección de los
funcionales angulares.


% FORMALIZACION_NUEVA de la elevacion de N39, compatible con (rho9,q9),
% T_d y Upd_t; arquitectura de residuo/cociente y orientacion preexistente.
% Fragmento de investigacion separado. No selecciona las reglas CKM N42.
\subsubsection{Transporte de acarreo y holonomía entera del cruce aritmético}
\label{subsec:transporte-entero-cruce}

La lectura modular del cruce admite una realización integral que mantiene
la información acumulada al atravesar el borde. Se desarrolla aquí esa
realización mediante los cocientes de transporte. Su dominio es una coordenatización
elevada del soporte APP; la posterior asignación de una región a un sector
angular conserva su propio mapa de incidencia.

\paragraph{Coordenatización nonádica y transporte elevado.}
Se emplea la descomposición exacta
\[
 \rho_9(n)=1+((n-1)\bmod9),\qquad
 q_9(n)=\left\lfloor\frac{n-1}{9}\right\rfloor,
 \qquad n=\rho_9(n)+9q_9(n).
\]
Una posición elevada es
$z=(x,y,q_x,q_y)\in D_9^2\times\mathbb Z^2$, donde
$X=x+9q_x$ e $Y=y+9q_y$. El incremento entero $u$ de la primera
coordenada actúa por
\[
 T^x_u(x,y,q_x,q_y)=
 \bigl(\rho_9(x+u),y,q_x+c_9(x;u),q_y\bigr),\qquad
 c_9(x;u)=\left\lfloor\frac{x+u-1}{9}\right\rfloor.
\]
La segunda coordenada se transporta de la misma manera. Los avances
cardinales corresponden a $u=\pm1$; el incremento nulo conserva el estado.
El cociclo satisface
\[
 c_9(x;u+v)=c_9(x;u)+c_9(\rho_9(x+u);v).
\]
En efecto, ambas partes son el cociente de la única descomposición de
$x+u+v$ con residuo en $D_9$. Por consiguiente,
$T^x_vT^x_u=T^x_{u+v}$ y $T^x_{-u}=(T^x_u)^{-1}$; los transportes
de las dos coordenadas conmutan. El estado conserva el número de cruces
de borde, incluso cuando la posición visible retorna.

\paragraph{Hojas elevadas y actualización de los registros.}
La realización integral se fija prolongando las operaciones aritméticas
APP a las coordenadas transportadas:
\begin{align*}
 \widetilde S(z)&=X+Y=x+y+9(q_x+q_y),\\
 \widetilde P(z)&=XY
 =xy+9(q_xy+q_yx)+81q_xq_y.
\end{align*}
En cada posición se conserva el registro de ambas hojas
\[
 \bigl(\rho_9(\widetilde S),q_9(\widetilde S),
       \rho_9(\widetilde P),q_9(\widetilde P)\bigr).
\]
Si $e:z\to z'$ es un enlace y $T$ la hoja seleccionada, su incremento
íntegro se obtiene del registro antes y después del transporte:
\begin{equation}
 \delta_e\widetilde T=
 \rho_9(\widetilde T(z'))-\rho_9(\widetilde T(z))
 +9\bigl(q_9(\widetilde T(z'))-q_9(\widetilde T(z))\bigr).
 \label{eq:cruce-incremento-registro}
\end{equation}
Esta identidad explicita la conservación: el salto del residuo y el
incremento del cociente reconstruyen juntos la variación de la hoja.
La selección de hoja, el transporte cardinal y la actualización del
registro son operaciones distintas. En esta coordenatización se compone primero la
selección ordenada de hojas $(T_x,T_y)$ con $T_d$ y con la actualización
de sus cocientes; la fórmula no postula que todo el TPK quede reducido
a esta representación local.

\begin{proposition}[Curvatura integral y transporte de frontera]
\label{prop:cruce-stokes-integral}
Defínanse
$\widetilde A_x=\delta_x\widetilde T_x$ y
$\widetilde A_y=\delta_y\widetilde T_y$, mediante
\eqref{eq:cruce-incremento-registro}, y sea
\[
 \widetilde F_{xy}(X,Y)=
 \widetilde A_x(X,Y)+\widetilde A_y(X+1,Y)
 -\widetilde A_x(X,Y+1)-\widetilde A_y(X,Y).
\]
Las elecciones homogéneas tienen curvatura cero. Las elecciones mixtas
$(\widetilde S,\widetilde P)$ y
$(\widetilde P,\widetilde S)$ tienen curvatura, respectivamente,
$+1$ y $-1$ como enteros. Para el borde positivo de un rectángulo
$R_{a,b}$ de lados enteros $a,b>0$,
\[
 \widetilde W_{SP}(\partial R_{a,b})=ab,\qquad
 \widetilde W_{PS}(\partial R_{a,b})=-ab.
\]
El resultado es independiente del punto base y del número de cruces
de las costuras nonádicas.
\end{proposition}

\begin{proof}
La reconstrucción exacta de ambas hojas da
\[
 \delta_x\widetilde S=\delta_y\widetilde S=1,
 \qquad\delta_x\widetilde P=Y,
 \qquad\delta_y\widetilde P=X.
\]
Por tanto, la conexión $SP$ es $(1,X)$, con curvatura
$1+(X+1)-1-X=1$; la conexión $PS$ es $(Y,1)$, con curvatura
$Y+1-(Y+1)-1=-1$. Para una hoja única se cancelan las diferencias
mixtas. La suma de curvaturas de las $ab$ plaquetas cancela cada
enlace interior con su inverso y conserva los enlaces de frontera.
Así se obtiene $\pm ab$ en $\mathbb Z$. Los cocientes transportados
garantizan que la misma identidad rige al atravesar cualquier costura.
\end{proof}

El acumulador de frontera puede conservarse a su vez en forma nonádica:
$W=r_W+9q_W$. Para una contribución entera $w_e$, se actualiza mediante
\[
 r_W'=\rho_9(r_W+w_e),\qquad
 q_W'=q_W+\left\lfloor\frac{r_W+w_e-1}{9}\right\rfloor.
\]
La composición por enlaces proporciona la misma suma que su evaluación
entera directa. En esta coordenatización, el cero entero tiene coordenadas
$(r_W,q_W)=(9,-1)$; esta es una representación aritmética del acumulador,
sin identificación adicional con los distintos ceros orientados del TRIT.

\paragraph{Tres operaciones sobre la orientación.}
Sea $\gamma$ una ruta cardinal cerrada. Su inversión temporal
$\gamma^{-1}$ invierte el orden de enlaces y la orientación de cada uno;
por aditividad,
\[
 \widetilde W_{SP}(\gamma^{-1})=-\widetilde W_{SP}(\gamma).
\]
El intercambio de las hojas satisface asimismo
\[
 \widetilde W_{PS}(\gamma)=-\widetilde W_{SP}(\gamma).
\]
Para comprobarlo en cualquier ruta, obsérvese que la suma de ambas
conexiones es el incremento exacto de
$\widetilde S+\widetilde P=X+Y+XY$; su integral sobre un lazo es cero.

La semivuelta simultánea de las dos direcciones actúa, por otra parte,
mediante $J_c(X,Y)=(2c_x-X,2c_y-Y)$, con $2c\in\mathbb Z^2$,
de modo que preserva la coordenatización reticular. Para un lazo se cumple
\[
 \widetilde W_{SP}(J_c\gamma)=\widetilde W_{SP}(\gamma).
\]
En efecto, la contribución horizontal cerrada es cero y la vertical
se transforma como
$\sum (2c_x-X)(-\delta Y)=\sum X\delta Y$, pues
$\sum\delta Y=0$. Una semivuelta espacial preserva, por tanto, la
orientación areal de esta coordenatización. Su combinación con la inversión temporal
o con el intercambio de hojas se calcula componiendo las operaciones
correspondientes.

La diferencia es pertinente para el registro electrónico: la inversión
simultánea de los dos desplazamientos cada $27$ avances produce, en
ventanas de $9$, la sucesión
\[
 \tau_j=(-1)^{\lfloor j/3\rfloor},\qquad 0\leq j<12,
 \qquad (+,+,+,-,-,-,+,+,+,-,-,-).
\]
Esta sucesión registra orientaciones de avance. La holonomía requiere
además los enlaces recorridos, las hojas seleccionadas y los cocientes
acumulados. Así, el transporte de la orientación electrónica se conserva
sin identificar una semivuelta de ambos ejes con la inversión temporal
de un lazo.

\paragraph{Aplicación al cruce marcado y alcance sectorial.}
Para el cruce rectangular de lados $4$ y $7$, esta realización evalúa
exactamente
\[
 \widetilde W_{SP}=28=1+9\cdot3,\qquad
 \widetilde W_{SP}(\gamma^{-1})=-28=8+9(-4).
\]
El entero conserva el área orientada; el residuo aislado representa sólo
su clase nonádica. Por ejemplo, desde $(7,1)$ los representantes
periódicos de los enlaces suman $-35$, mientras que la corrección de
acarreo añade $63$ y recupera $28$.

En la incidencia marcada del lector angular, el $4$ es
$|H_e\cap P_\pi|$ y el $7$ es el pivote marcado $p_\varphi$.
La lectura rectangular aporta un realizador entero del producto
$|H_e\cap P_\pi|\,p_\varphi$ una vez declarado ese cruce. El
mapa que identifique esa región con el sector $12$ debe conservar
la marca y la orientación del sistema sectorial de
la Definición~\ref{def:ii09-sistema-sectorial}. Las asignaciones
sectoriales $12,23,13$ y sus combinaciones de $A,h,\Delta^\circ$
son una especificación adicional a la ley de holonomía rectangular.
El presente resultado materializa la elevación de área, el acarreo y
sus involuciones, con esos datos de partida explícitos.


\begin{definition}[Sistema sectorial de incidencia angular]
\label{def:ii09-sistema-sectorial}
El sistema sectorial consta del sexteto ordenado de invariantes
\[
 (b,p,q,h_6,o_8,t)=(4,5,7,6,8,3)
\]
y de las aplicaciones lineales, sobre coordenadas angulares
\((a,u,d)\),
\begin{align}
 \mathscr A_{12}^{\rm ang}(a,u,d)&=2a-pu+bq\,d,\label{eq:ii09-sector12}\\
 \mathscr A_{23}^{\rm ang}(a,u,d)&=qu-\frac{p^2}{2}\,d,\label{eq:ii09-sector23}\\
 \mathscr A_{13}^{\rm ang}(a,u,d)&=a-pb\,u-\frac{o_8-h_6}{t}\,d.
 \label{eq:ii09-sector13}
\end{align}
Los símbolos \(h_6,o_8\) conservan los cardinales de la incidencia
marcada; \(t=3\) es la normalización trítica declarada de ese sector.
El sexteto y las tres reglas forman conjuntamente el sistema de
partida de esta construcción.
\end{definition}

\begin{theorem}[Evaluación exacta del sistema sectorial]
\label{thm:ii09-matriz-sectorial}
La aplicación conjunta
\((\mathscr A_{12}^{\rm ang},\mathscr A_{23}^{\rm ang},
\mathscr A_{13}^{\rm ang})^{\mathsf T}\) de la
Definición~\ref{def:ii09-sistema-sectorial} tiene por matriz única
\[
 \mathsf M_{\rm CKM}=
 \begin{pmatrix}
 2&-5&28\\
 0&7&-25/2\\
 1&-20&-2/3
 \end{pmatrix}.
\]
Las tres columnas son las imágenes de los vectores coordenados
\((1,0,0)\), \((0,1,0)\) y \((0,0,1)\) bajo esas aplicaciones.
\end{theorem}
\begin{proof}
Las operaciones sectoriales dan
\[
 bq=4\cdot7=28,\qquad p^2/2=5^2/2=25/2,
 \qquad pb=5\cdot4=20,\qquad
 (o_8-h_6)/t=(8-6)/3=2/3.
\]
Por tanto, las imágenes sucesivas de la base coordenada son
\[
 c_A=(2,0,1)^{\mathsf T},\qquad
 c_h=(-5,7,-20)^{\mathsf T},\qquad
 c_\Delta=(28,-25/2,-2/3)^{\mathsf T}.
\]
Una aplicación lineal queda determinada por sus imágenes sobre una
base. La matriz cuyas columnas son estas imágenes es la indicada.
La prueba utiliza el sistema sectorial completo de la definición.
\end{proof}

Sobre las coordenadas previamente producidas se toma
\[
 (a,u,d)=\left(A_{\deg},\frac{C^*_{\deg}}6,
                       \frac{180}{\pi}\Delta_4\right).
\]
Así quedan individualizadas las contribuciones directa, conjugada
y torsional a cada ángulo. En particular, la diferencia
\(o_8-h_6=2\) se conserva en el coeficiente del defecto torsional
del sector \(13\). Esta presencia concreta permite seguir la
incidencia hasta el lector angular, manteniendo separados el
sistema sectorial, su evaluación y la parametrización unitaria.


\subsection{Coordenatización racional CKM y unidades angulares}

Escribimos \(A_{\rm deg}\) y \(C^*_{\rm deg}\) para las coordenadas en
grados de los ángulos HMT ya generados. Se utilizan las abreviaturas
\begin{equation}
 h_\theta=\frac{C^*_{\rm deg}}6,
 \qquad
 \Delta_{\rm deg}=\frac{180}{\pi_{\rm HMT}}\Delta_4,
 \qquad
 x_\theta=(A_{\rm deg},h_\theta,\Delta_{\rm deg})^{\mathsf T}.
 \label{eq:ii-angular-chart}
\end{equation}
Aquí \(h_\theta\) es una coordenada angular, distinta de la constante de
Planck \(h\). La conversión a grados es un cambio de coordenadas posterior; no
introduce un valor convencional de \(\pi\) en el generador.

\begin{definition}[Lector angular quark]
El lector heredado \(\mathsf M_{\rm CKM}:\mathbb R^3\to\mathbb R^3\)
evalúa las coordenadas
\begin{equation}
 \boldsymbol\theta_{\rm deg}=\mathsf M_{\rm CKM}x_\theta,
 \qquad
 \mathsf M_{\rm CKM}=
 \begin{pmatrix}
 2&-5&28\\
 0&7&-25/2\\
 1&-20&-2/3
 \end{pmatrix},
 \qquad
 \delta_{\rm deg}=9A_{\rm deg}.
 \label{eq:ii-ckm-map}
\end{equation}
El orden de las componentes es \((12,23,13)\). Las funciones
trigonométricas posteriores se evalúan en
\(\vartheta_{ab}=(\pi_{\rm HMT}/180)\theta_{ab,\rm deg}\) y
\(\delta=(\pi_{\rm HMT}/180)\delta_{\rm deg}\).
\end{definition}

\begin{proposition}[Reconstrucción racional de las coordenadas]
El lector tiene determinante \(-3857/6\) e inversa
\begin{equation}
 \mathsf M_{\rm CKM}^{-1}=
 \begin{pmatrix}
 1528/3857&3380/3857&801/3857\\
 75/3857&176/3857&-150/3857\\
 6/551&-30/551&-12/551
 \end{pmatrix}.
 \label{eq:ii-ckm-inverse}
\end{equation}
En particular,
\begin{align}
 A_{\rm deg}
 &=\frac{1528\theta_{12,\rm deg}+3380\theta_{23,\rm deg}
          +801\theta_{13,\rm deg}}{3857},\nonumber\\
 C^*_{\rm deg}
 &=\frac{6(75\theta_{12,\rm deg}+176\theta_{23,\rm deg}
          -150\theta_{13,\rm deg})}{3857},\nonumber\\
 \Delta_{\rm deg}
 &=\frac{6\theta_{12,\rm deg}-30\theta_{23,\rm deg}
          -12\theta_{13,\rm deg}}{551}.
 \label{eq:ii-ckm-recover}
\end{align}
La fase satisface el control lineal exacto
\begin{equation}
 3857\delta_{\rm deg}
 =9(1528\theta_{12,\rm deg}+3380\theta_{23,\rm deg}
      +801\theta_{13,\rm deg}).
 \label{eq:ii-ckm-phase-check}
\end{equation}
\end{proposition}

\begin{proof}
La expansión del determinante de \eqref{eq:ii-ckm-map} da
\(-3857/6\). La multiplicación de la matriz indicada en
\eqref{eq:ii-ckm-inverse} por \(\mathsf M_{\rm CKM}\), en ambos
órdenes, da \(I_3\). Sus filas recuperan las tres componentes de
\(x_\theta\); la segunda se multiplica por seis para recuperar
\(C^*_{\rm deg}\). Finalmente se usa
\(\delta_{\rm deg}=9A_{\rm deg}\).
\end{proof}

La inversión reconstruye coordenadas del lector una vez fijado éste. No
convierte los ángulos observados en antecedentes de \(A\), \(C^*\) o
\(\Delta_4\).

\subsection{Descomposición geométrica bajo inversión de hoja}
\label{sec:ii09-descomposicion-paridad}

En la coordenatización angular, la inversión de hoja conserva \(a,d\) e
invierte \(u\). El vector expresado admite la descomposición
\[
 \boldsymbol\theta_+=a c_A+u c_h+d c_\Delta,
 \qquad
 \boldsymbol\theta_-=a c_A-u c_h+d c_\Delta.
\]
\begin{proposition}[Componentes par e impar del transporte angular]
\label{prop:ii09-paridad-angular}
Los dos vectores satisfacen
\[
 \frac{\boldsymbol\theta_++\boldsymbol\theta_-}{2}
       =a c_A+d c_\Delta,\qquad
 \frac{\boldsymbol\theta_+-\boldsymbol\theta_-}{2}=u c_h.
\]
La contribución que cambia de signo ocupa la recta
\(\operatorname{span}(c_h)\); la componente conservada ocupa
\(\operatorname{span}(c_A,c_\Delta)\).
\end{proposition}
\begin{proof}
La suma cancela los términos \(u c_h\), y la diferencia cancela
los otros dos. Las columnas son linealmente independientes, pues
\(\det\mathsf M_{\rm CKM}=-3857/6\). Por ello las dos
componentes forman una descomposición directa.
\end{proof}

La misma inversión actúa sobre la evaluación de Barbero--Immirzi:
intercambia \(q_+\) y \(q_-\), e invierte la componente bilineal
en \(A\) y \(C^*\). En consecuencia, el funcional de incidencia
es impar. Un observable simétrico en ambos canales conserva su
valor. Esta diferencia de paridad expresa cómo una misma
operación sobre la estructura produce respuestas distintas según
el lector: orientación areal, información simétrica y vector de
mezcla conservan información diferente del mismo antecedente.


\subsection{Transporte de la inversión de hoja}

La conjugación de las coordenatizaciones angulares mantiene \(A_{\rm deg}\) y
\(\Delta_{\rm deg}\), e invierte \(h_\theta\). Definimos
\begin{equation}
 D_\theta=\operatorname{diag}(1,-1,1),\qquad
 c_h=(-5,7,-20)^{\mathsf T},\qquad
 \ell_h=\frac1{3857}(75,176,-150).
 \label{eq:ii-ckm-sheet-data}
\end{equation}
Si \(c_A=(2,0,1)^{\mathsf T}\) y
\(c_\Delta=(28,-25/2,-2/3)^{\mathsf T}\), entonces
\(\ell_hc_A=\ell_hc_\Delta=0\) y \(\ell_hc_h=1\).

\begin{theorem}[Reflexión racional y métrica transportada]
El operador
\begin{equation}
 \mathcal R_{\rm CKM}
 =\mathsf M_{\rm CKM}D_\theta\mathsf M_{\rm CKM}^{-1}
 =I_3-2c_h\ell_h
 \label{eq:ii-ckm-reflection}
\end{equation}
satisface
\begin{equation}
 \mathcal R_{\rm CKM}^2=I_3,\qquad
 \det\mathcal R_{\rm CKM}=-1,\qquad
 \mathcal R_{\rm CKM}\mathsf M_{\rm CKM}
 =\mathsf M_{\rm CKM}D_\theta.
 \label{eq:ii-ckm-intertwining}
\end{equation}
Es una reflexión ortogonal para la métrica positiva
\begin{equation}
 G_\theta=(\mathsf M_{\rm CKM}^{-1})^{\mathsf T}
                 \mathsf M_{\rm CKM}^{-1},\qquad
 \mathcal R_{\rm CKM}^{\mathsf T}G_\theta\mathcal R_{\rm CKM}=G_\theta.
 \label{eq:ii-ckm-metric}
\end{equation}
\end{theorem}

\begin{proof}
El operador \(P_h=c_h\ell_h\) verifica \(P_h^2=P_h\), tiene imagen
\(\mathbb Rc_h\) y núcleo \(\operatorname{span}(c_A,c_\Delta)\).
Por ello \((I_3-2P_h)^2=I_3\), con autovalores \(1,1,-1\).
La acción sobre esas tres columnas es la de \(D_\theta\) en la base
original, lo que prueba la conjugación y el entrelazamiento.
Para \(v\ne0\),
\(v^{\mathsf T}G_\theta v=\|\mathsf M_{\rm CKM}^{-1}v\|^2>0\).
La última identidad resulta de \(D_\theta^{\mathsf T}D_\theta=I_3\).
\end{proof}

Su expresión completa, conservada del soporte anterior, es
\begin{equation}
 \mathcal R_{\rm CKM}=
 \begin{pmatrix}
 4607/3857&1760/3857&-1500/3857\\
 -150/551&199/551&300/551\\
 3000/3857&7040/3857&-2143/3857
 \end{pmatrix}.
 \label{eq:ii-ckm-reflection-matrix}
\end{equation}
La fase \(\delta_{\rm deg}=9A_{\rm deg}\) permanece fija bajo esta
reflexión. En consecuencia, el transporte de hoja no debe identificarse
con la conjugación CP de amplitudes, que cambia el signo de la fase compleja.
La paridad combinatoria de una ruta, su conjugación de hoja y su
representación física mantienen los dominios diferenciados en
\cite{cpt,mezcla}.

\subsection{Paridad del registro y realización CP}

Para conservar el soporte discreto de esa distinción, sea \(\rho\) una
ruta registrada de longitud \(L\), con direcciones
\(d_t\in\{N,E,S,W\}\) y hojas \(\ell_t\in\{\Sigma,\Pi\}\).
Los estados intermedios son los producidos por el transporte TPK.
La reflexión \(P\) actúa por \((i,j)\mapsto(i,-j)\pmod9\) y
permuta \(E,W\); \(C_{\rm hoja}\) permuta \(\Sigma,\Pi\);
\(T_{\rm reg}\) invierte el orden del registro y reemplaza cada
dirección por su opuesta. La inversión se efectúa sobre el registro
enriquecido, no sobre una proyección que haya olvidado su predecesor.

\begin{proposition}[Invariancia de la acción combinatoria de ruta]
Para \(\lambda\ge0\), el funcional adimensional
\begin{equation}
 \mathcal A_\lambda(\rho)=
 \sum_{t=2}^{L}\left[1-\delta_{d_t,d_{t-1}}
             +\lambda(1-\delta_{\ell_t,\ell_{t-1}})\right]
 \label{eq:ii-discrete-route-action}
\end{equation}
es invariante bajo \(P\), \(C_{\rm hoja}\), \(T_{\rm reg}\) y
sus composiciones.
\end{proposition}

\begin{proof}
Una permutación global del alfabeto conserva la igualdad entre vecinos.
La inversión del orden preserva el número de cambios consecutivos.
Por tanto, cada operación conserva por separado los giros de dirección
y los cambios de hoja que suma \eqref{eq:ii-discrete-route-action}.
\end{proof}

\(\lambda\) es el parámetro declarado de esta colección de funcionales;
\(\mathcal A_\lambda\) no se identifica con una sección dimensional de
Planck. La invariancia del contador tampoco implica que cada lector
orientado sea par. En una realización de mezcla \(\mathcal V\), la
compatibilidad CP toma la forma precisa
\begin{equation}
 \mathcal V(CP_{\rm HMT}\rho)
 =D_u\,\overline{\mathcal V(\rho)}\,D_d^\dagger,
 \label{eq:ii-cp-square}
\end{equation}
con los cambios diagonales de fase y los espacios de representación
declarados. El contador de ruta, el transporte racional de hoja y la
conjugación compleja son tres mapas distintos del soporte preservado.

\begin{remark}[Compatibilidad CP del lector físico]
La ecuación \eqref{eq:ii-cp-square} es la condición de
compatibilidad CP del lector físico completo.
Las demostraciones de esta sección establecen la invariancia del
registro y las identidades del lector complejo, pero no convierten
automáticamente \(\mathcal R_{\rm CKM}\) en \(CP_{\rm HMT}\).
Esta reserva se refiere a ese cuadrado de representación, no a la
existencia del transporte interno de rutas ya construido.
\end{remark}

\subsection{Amplitudes unitarias, fase CP e invariante de Jarlskog}

Escribimos \(c_{ab}=\cos\vartheta_{ab}\) y
\(s_{ab}=\sin\vartheta_{ab}\). La realización compleja del lector es
\begingroup
\small
\setlength{\arraycolsep}{3pt}
\begin{equation}
 V_{\rm CKM}=
 \begin{pmatrix}
 c_{12}c_{13}&s_{12}c_{13}&s_{13}e^{-i\delta}\\
 -s_{12}c_{23}-c_{12}s_{23}s_{13}e^{i\delta}&
 c_{12}c_{23}-s_{12}s_{23}s_{13}e^{i\delta}&s_{23}c_{13}\\
 s_{12}s_{23}-c_{12}c_{23}s_{13}e^{i\delta}&
 -c_{12}s_{23}-s_{12}c_{23}s_{13}e^{i\delta}&c_{23}c_{13}
 \end{pmatrix}.
 \label{eq:ii-ckm-unitary}
\end{equation}
\endgroup
La estructura compleja utilizada es la realización de la estructura
orientada ya construida; la parametrización unitaria representa el resultado
angular, no determina sus coeficientes.

\begin{proposition}[Unitariedad e invariante de fase]
La matriz \eqref{eq:ii-ckm-unitary} es unitaria. El cuarteto
\begin{equation}
 J_{\rm CKM}
 =\operatorname{Im}
   (V_{11}V_{22}\overline{V_{12}}\,\overline{V_{21}})
 =c_{12}c_{23}c_{13}^2s_{12}s_{23}s_{13}\sin\delta
 \label{eq:ii-jarlskog}
\end{equation}
es invariante bajo \(V\mapsto D_uVD_d\), con \(D_u,D_d\) diagonales
unitarias, y cambia de signo bajo \(V\mapsto\overline V\).
\end{proposition}

\begin{proof}
La matriz es el producto \(R_{23}U_{13}(\delta)R_{12}\), donde
\begin{equation}
\begin{aligned}
 R_{12}&=\begin{pmatrix}c_{12}&s_{12}&0\\-s_{12}&c_{12}&0\\0&0&1\end{pmatrix},
 \\
 U_{13}(\delta)&=
 \begin{pmatrix}c_{13}&0&s_{13}e^{-i\delta}\\0&1&0\\
 -s_{13}e^{i\delta}&0&c_{13}\end{pmatrix},
 \\
 R_{23}&=\begin{pmatrix}1&0&0\\0&c_{23}&s_{23}\\0&-s_{23}&c_{23}\end{pmatrix}.
\end{aligned}
 \label{eq:ii-ckm-factorization}
\end{equation}
Cada factor es unitario por \(c_{ab}^2+s_{ab}^2=1\).
Al sustituir las cuatro entradas del cuarteto, su parte imaginaria es
\(c_{12}s_{12}c_{13}^2c_{23}s_{23}s_{13}
(c_{12}^2+s_{12}^2)\sin\delta\), lo que da
\eqref{eq:ii-jarlskog}. Los factores de fase de cada fila y cada columna
se cancelan con sus conjugados. La conjugación compleja invierte la parte
imaginaria.
\end{proof}

Un \(J_{\rm CKM}\) no nulo obstruye hacer real toda la matriz mediante
esos cambios de fase. Su interpretación física mantiene además el sector
y la representación correspondientes; no es un corrector escalar añadido
a las masas. Si los espectros presentan degeneraciones, sus libertades de
cambio de base deben conservarse antes de formular una conclusión CP.

\subsection{Conjugación de amplitudes e invariante orientado}
\label{sec:ii09-conjugacion-amplitudes}

La realización compleja conserva los tres ángulos de mezcla y
la fase de interferencia como coordenadas explícitas. Escribamos
\(V(\boldsymbol\vartheta,\delta)
=R_{23}U_{13}(\delta)R_{12}\), con las rotaciones y la fase
definidas en el desarrollo unitario. Los argumentos de las
funciones trigonométricas son sus representantes en radianes.

\begin{proposition}[Conjugación en la realización de mezcla]
\label{prop:ii09-conjugacion-unitaria}
Para los tres ángulos reales y la fase real declarados,
\[
 V(\boldsymbol\vartheta,-\delta)
       =\overline{V(\boldsymbol\vartheta,\delta)},
 \qquad
 J(\boldsymbol\vartheta,-\delta)
       =-J(\boldsymbol\vartheta,\delta).
\]
El invariante \(J\) conserva además su valor bajo los cambios
diagonales de fase de las bases cargadas.
\end{proposition}
\begin{proof}
Las matrices \(R_{12}\) y \(R_{23}\) tienen entradas reales.
En \(U_{13}\), el cambio \(\delta\mapsto-\delta\) conjuga
los factores exponenciales. La conjugación del producto da la
primera identidad. En el cuarteto
\(V_{11}V_{22}\overline{V_{12}}\overline{V_{21}}\), los
factores diagonales de cada fila y columna se cancelan. La
conjugación invierte el signo de su parte imaginaria. Se obtiene
equivalentemente el resultado sustituyendo en
\[
 J=c_{12}c_{23}c_{13}^{2}s_{12}s_{23}s_{13}\sin\delta.
\]
\end{proof}

En el sector de tres generaciones, el valor interno no nulo de
\(J\) expresa la obstrucción a realizar todas las amplitudes
mediante una matriz real por cambios diagonales de fase. Esta
es la lectura de violación de carga--paridad de la matriz
construida. La lectura de rutas conserva, a su vez, la operación
discreta de conjugación y paridad y su representación. Su
compatibilidad completa se expresa mediante el cuadrado
\[
 \mathcal V(CP_{\rm HMT}\rho)
  =D_u\,\overline{\mathcal V(\rho)}\,D_d^{\dagger},
\]
con los espacios de representación y los cambios de fase
declarados. La identidad de conjugación probada aquí fija
exactamente la acción que debe transportar ese cuadrado.


\Needspace{14\baselineskip}
\begin{samepage}
\begin{proposition}[Mezcla y espectro]
Sean \(B_u=\operatorname{diag}(a_1,a_2,a_3)\) y
\(B_d=\operatorname{diag}(b_1,b_2,b_3)\) los lectores hermíticos de masa
ya producidos en sus fibras. El transporte
\begin{equation}
 B_d^{(u)}=V_{\rm CKM}B_dV_{\rm CKM}^{\dagger}
 \label{eq:ii-ckm-mass-transport}
\end{equation}
conserva el espectro de \(B_d\), y satisface
\begin{equation}
 \left|\det[B_u,B_d^{(u)}]\right|
 =2|J_{\rm CKM}|
   \prod_{i<j}|a_i-a_j|\prod_{r<s}|b_r-b_s|.
 \label{eq:ii-ckm-commutator}
\end{equation}
\end{proposition}
\end{samepage}

\begin{proof}
La conjugación unitaria conserva el polinomio característico. Si
\(B=B_d^{(u)}\), las entradas del conmutador son
\([B_u,B]_{ij}=(a_i-a_j)B_{ij}\), con diagonal nula. Su determinante es
\[
 2i(a_1-a_2)(a_2-a_3)(a_3-a_1)
     \operatorname{Im}(B_{12}B_{23}B_{31}).
\]
Al sustituir \(B_{ij}=\sum_r b_rV_{ir}\overline{V_{jr}}\), la parte
imaginaria es un polinomio alternante de grado tres en los \(b_r\).
Es divisible por sus tres diferencias: si dos \(b_r\) coinciden,
\(B=bI+(b'-b)vv^\dagger\) y el producto fuera de diagonal es real.
La expansión de un coeficiente, usando la ortogonalidad de las columnas,
da el cuarteto de \eqref{eq:ii-jarlskog}, con el signo del orden elegido.
Tomar valores absolutos elimina esa convención y produce
\eqref{eq:ii-ckm-commutator}, incluida su anulación en las degeneraciones.
\end{proof}

\begin{remark}[Evaluación y comparación]
Las coordenadas decimales CKM de la
tabla~\ref{tab:ii-ckm-comparacion} son evaluaciones posteriores del
lector anterior. La sección de comparación cuantitativa presenta también
los ángulos y los módulos de las amplitudes. Esas evaluaciones no se
utilizan para escoger sus columnas ni para demostrar
\eqref{eq:ii-ckm-intertwining} o \eqref{eq:ii-jarlskog}.
Este bloque no asigna significación metrológica nueva a esas cifras:
una comparación conjunta requiere su edición externa, definición de
observables y covarianza. Las desviaciones marginales no la sustituyen.
\end{remark}

\subsection{Registro leptónico, operadores y marcos espectrales}

El sector leptónico conserva otro dominio. La construcción espectral de
\eqref{eq:ii-neutrino-compression} reúne cuatro profundidades neutras de dimensiones
\(2,4,6,8\), antes de la compresión tridimensional:
\begin{equation}
 \mathcal H_\nu^{\rm reg}
 =\bigoplus_{j=1}^{4}\mathcal H_{\nu,G_j},\qquad
 \dim\mathcal H_\nu^{\rm reg}=20,\qquad
 M_\nu=\left.(P_\nu\mathcal M_\nu P_\nu)\right|_{\operatorname{Im}P_\nu},
 \quad\operatorname{rank}P_\nu=3.
 \label{eq:ii-neutrino-compression}
\end{equation}
Las cuatro profundidades no son cuatro sabores ni cuatro masas.
\(P_\nu\) es el proyector neutral heredado del lector de registro, y
\(M_\ell=\left.(P_\ell\mathcal M_\ell P_\ell)\right|_{\operatorname{Im}P_\ell}\),
con \(\operatorname{rank}P_\ell=3\), es el operador cargado comprimido.
La escala de esos operadores procede de la rama acción--reloj--masa;
ningún ángulo PMNS la determina.

\begin{definition}[Transporte de marcos leptónicos]
En la realización común tridimensional declarada, sean
\(F_\ell,F_\nu:\mathbb C^3\to\mathcal H_3\) marcos espectrales
ortonormales completos de los operadores cargado y neutro. Se fija la
identificación con \(\mathcal H_3\) antes de tomar su producto. El lector es
\begin{equation}
 U_{\rm PMNS}=F_\ell^\dagger F_\nu:
 \mathbb C^3_\nu\longrightarrow\mathbb C^3_\ell.
 \label{eq:ii-pmns-frames}
\end{equation}
Los marcos y la identificación proceden de la representación y del
registro; no se eligen a partir de una matriz experimental que se desea
reproducir.
\end{definition}

\begin{theorem}[Unitariedad, degeneraciones y naturalidad]
Bajo la identificación precedente, \(U_{\rm PMNS}\) es unitaria.
Si los espectros son simples, sus cambios de marco son fases diagonales
y permutaciones; fijado el orden espectral, permanecen las fases.
Si hay multiplicidades \(m_{\ell,\lambda}\) y \(m_{\nu,\eta}\), el
objeto independiente de esas bases es la clase doble
\begin{equation}
 \left(\prod_\lambda U(m_{\ell,\lambda})\right)
 \backslash U(3) /
 \left(\prod_\eta U(m_{\nu,\eta})\right).
 \label{eq:ii-pmns-double-class}
\end{equation}
Un refinamiento unitario común \(W\) que entrelace ambos operadores
conserva el transporte cuando los marcos se prolongan como
\(F'_\ell=WF_\ell\), \(F'_\nu=WF_\nu\).
\end{theorem}

\begin{proof}
La completitud en el mismo espacio da
\(F_\ell F_\ell^\dagger=F_\nu F_\nu^\dagger=I_{\mathcal H_3}\).
Por tanto,
\(U_{\rm PMNS}^\dagger U_{\rm PMNS}
=F_\nu^\dagger F_\ell F_\ell^\dagger F_\nu=I_3\), y análogamente
en el otro orden. Los cambios de base que preservan un operador hermítico
actúan por bloques dentro de sus espacios propios. Al sustituir
\(F_\ell D_\ell,F_\nu D_\nu\), resulta
\(U\mapsto D_\ell^\dagger UD_\nu\), que determina la clase doble.
Finalmente, \((WF_\ell)^\dagger WF_\nu=F_\ell^\dagger F_\nu\).
\end{proof}

La condición de espacio común es material. Para dos marcos isométricos en
un ambiente mayor, la ortonormalidad por separado sólo produce una
compresión; no autoriza cancelar \(F_\ell F_\ell^\dagger\) como si
fuera la identidad del ambiente.

\subsection{Núcleo de registro y suficiencia de rango}

Los lectores \(\Xi_r\) conservan los componentes de acarreo, hoja,
orientación, vacancia y holonomía del registro. Sobre el conjunto finito
de rutas admitidas se forma
\begin{equation}
 K_{\rm reg}(c,d)=
 \sum_{\substack{r\in\{\mathrm{car},\mathrm{hoj},\mathrm{ori},
                           \mathrm{vac},\mathrm{hol}\}\\ d_r>0}}
 \frac{\langle\Xi_r(c),\Xi_r(d)\rangle}{d_r},
 \qquad
 d_r=\sum_x\|\Xi_r(x)\|^2.
 \label{eq:ii-reg-kernel}
\end{equation}
La normalización procede de las normas del propio registro. Los
componentes de norma nula se omiten; no se reemplazan por términos tomados
de una matriz objetivo. Cada \(\Xi_r\) y cada clase de agregación debe
conservar su residencia de ruta y su compatibilidad con el refinamiento.

\begin{proposition}[Positividad y rango del registro]
La matriz \(K_{\rm reg}\) es positiva semidefinida. Si se agregan sus
vectores de características en tres sectores, la matriz de Gram
\(G_\nu\) resultante es invertible exactamente cuando esos tres vectores
agregados son linealmente independientes.
\end{proposition}

\begin{proof}
Se define
\(v_c=\bigoplus_{r:d_r>0}d_r^{-1/2}\Xi_r(c)\).
Entonces \(K_{\rm reg}(c,d)=\langle v_c,v_d\rangle\), y para
cualesquiera coeficientes \(z_c\),
\(\sum_{c,d}\overline z_cK_{\rm reg}(c,d)z_d
=\|\sum_c z_cv_c\|^2\ge0\).
La misma identidad vale para los tres vectores agregados.
El núcleo de su matriz de Gram consiste exactamente en sus relaciones
lineales, de donde resulta el criterio de invertibilidad.
\end{proof}

Se distingue esa Gram de un sector de la matriz \emph{cruzada} entre las
clases cargadas \(L_a\) y neutras \(N_b\), expresada en \cite{mezcla}:
\begin{equation}
 G_{ab}=\frac1{\sqrt{|L_a||N_b|}}
       \sum_{c\in L_a}\sum_{d\in N_b}K_{\rm reg}(c,d).
 \label{eq:ii-cross-gram}
\end{equation}
Las clases no vacías proceden del clasificador interno. La matriz cruzada
no tiene por qué ser hermítica, a diferencia de una Gram construida con
los mismos tres vectores en ambos lados.

\begin{proposition}[Factor polar y su alcance]
Si \(G\) es invertible,
\begin{equation}
 U_G=G(G^\dagger G)^{-1/2}
 \label{eq:ii-polar}
\end{equation}
es unitaria. Si \(G\) es además hermítica positiva definida, entonces
\(U_G=I_3\). En consecuencia, el rango de una Gram positiva certifica
la suficiencia de los observables, pero por sí solo no determina una
mezcla no trivial entre los marcos cargado y neutro.
\end{proposition}

\begin{proof}
\(G^\dagger G\) es positiva definida; su raíz positiva es invertible y
\[
 U_G^\dagger U_G=(G^\dagger G)^{-1/2}(G^\dagger G)
 (G^\dagger G)^{-1/2}=I_3.
\]
Por dimensión finita, también
\(U_GU_G^\dagger=I_3\). Si \(G>0\), la raíz positiva de
\(G^\dagger G=G^2\) es \(G\), y \(U_G=I_3\).
\end{proof}

\begin{remark}[Factorización polar y marcos espectrales]
Las ecuaciones \eqref{eq:ii-reg-kernel} y \eqref{eq:ii-cross-gram}
conservan, respectivamente, el criterio de rango del registro y
la matriz cruzada de los sectores cargado y neutro. Esta exposición
separa sus funciones para no convertir una Gram positiva en un selector
artificial de ángulos. Identificar \(U_G\) con
\eqref{eq:ii-pmns-frames} requiere la compatibilidad del núcleo cruzado
con los marcos espectrales declarados. La identidad de un factor polar
o la mera unitariedad no sustituyen esa identificación. No se expresa
aquí una matriz numérica PMNS nueva.
\end{remark}

\subsection{Control de la fase de rango uno}

Sea \(b_K\) el testigo normalizado de incidencia de la fuente y
\(Q_K=|b_K\rangle\langle b_K|\). El operador
\begin{equation}
 \Phi_K=I+(e^{i\pi_{\rm HMT}/6}-1)Q_K
 \label{eq:ii-rank-one-phase}
\end{equation}
es unitario en su ambiente. Para marcos isométricos se considera
\(U_\Phi=F_\ell^\dagger\Phi_KF_\nu\).

\begin{proposition}[Criterio exacto para una compresión unitaria]
El operador \(U_\Phi\) es una contracción. Es unitario exactamente
cuando
\begin{equation}
 \Phi_K(\operatorname{Im}F_\nu)=\operatorname{Im}F_\ell.
 \label{eq:ii-phase-compatibility}
\end{equation}
\end{proposition}

\begin{proof}
Sean \(P_\ell=F_\ell F_\ell^\dagger\) y
\(P_\nu=F_\nu F_\nu^\dagger\). Entonces
\(U_\Phi^\dagger U_\Phi
=F_\nu^\dagger\Phi_K^\dagger P_\ell\Phi_KF_\nu\le I_3\).
La igualdad equivale a
\((I-P_\ell)\Phi_KF_\nu=0\), es decir, a la inclusión de la
imagen izquierda de \eqref{eq:ii-phase-compatibility} en la derecha.
Ambas tienen dimensión tres, de modo que la inclusión es igualdad.
\end{proof}

La fuente \cite{mezcla} conserva este ansatz mínimo como control
negativo de identificación física. Su unitariedad ambiental y aun el
cumplimiento de \eqref{eq:ii-phase-compatibility} no reemplazan el
registro completo ni los marcos de \eqref{eq:ii-pmns-frames}. El
control no se elimina al reunir la propagación angular en el artículo II.

\begingroup\fontsize{10}{11.7}\selectfont
\setlength{\parskip}{2pt}\setlength{\abovedisplayskip}{5pt}\setlength{\belowdisplayskip}{5pt}
\subsection{Fase leptónica y criterio Dirac--Majorana}

Los cambios de fase de bases deben distinguirse de los datos de
autorconjugación del sector neutro. Para cualquier transporte unitario
leptónico puede formarse el cuarteto
\begin{equation}
 J_\ell=\operatorname{Im}
 \bigl(U_{e1}U_{\mu2}\overline{U_{e2}}\,
                        \overline{U_{\mu1}}\bigr).
 \label{eq:ii-leptonic-quartet}
\end{equation}
La cancelación de fases demostrada para \eqref{eq:ii-jarlskog} se
aplica igualmente aquí. El cuarteto no fija todos los datos de la
realización neutra ni determina por sí solo su tipo Dirac o Majorana.

La realización de Majorana se especifica mediante la
autorconjugación antiunitaria \(\mathcal C_\nu\) compatible con masa y
dinámica:
\begin{equation}
 \mathcal C_\nu^2=I,\qquad
 \mathcal C_\nu M_\nu=M_\nu\mathcal C_\nu,\qquad
 \mathcal C_\nu H_\nu=H_\nu\mathcal C_\nu,
 \label{eq:ii-majorana-compatibility}
\end{equation}
con dominios preservados. La representación se realiza en la forma real
fija correspondiente. La alternativa Dirac conserva además una carga
\(Q\) que distingue el modo y su conjugado:
\begin{equation}
 \mathcal C_\nu Q\mathcal C_\nu^{-1}=-Q,\qquad Q\psi\ne0.
 \label{eq:ii-dirac-charge}
\end{equation}
Estos enunciados conservan sus datos representacionales. La neutralidad,
la involución de palabras dodecafásicas y un sector topológico de Ising
no deciden separadamente la alternativa.

\begin{remark}[Libertad de bases y fases de Majorana]
La clase doble \eqref{eq:ii-pmns-double-class} describe la libertad de
bases del par hermítico. Si se fija además una forma real de Majorana,
los cambios de fase admisibles deben preservar ese dato adicional;
no pueden trasladarse sin más todas las libertades de una representación
Dirac. El presente bloque conserva el criterio de autorconjugación y
no asigna valores a fases Majorana ni masas neutrínicas. Tales valores
no se obtienen de la neutralidad ni de una polarización de rango uno.
\end{remark}

\Needspace{8\baselineskip}
\subsection{Cierre del transporte angular}

La construcción del sector quark determina primero las coordenadas angulares, después
las amplitudes complejas y, finalmente, los invariantes de fase y el
transporte de operadores. La construcción leptónica determina operadores de
registro y marcos espectrales antes del cambio de base. Las dos rutas
conservan su residencia en el estado enriquecido común y en la estructura
discreta del continuo; ninguna introduce una tabla experimental como
selector de los coeficientes o de los estados.

Los controles de esta sección son locales y explícitos: inversión y
entrelazamiento racional CKM; conservación espectral por conjugación;
completitud y espacio común de los marcos PMNS; rango del registro;
compatibilidad de la matriz cruzada; y preservación de los datos de
autorconjugación. Así se incorpora la propagación angular a la teoría de
acción, área e información sin hacer de la mezcla una segunda ley de
masas ni convertir el lector de Barbero en su causa.

\par\endgroup
